% !TEX root = ../main.tex

\chapter{Model Matematis dan Metode}
\label{chap:metode}

Bab ini menyajikan derivasi matematis yang menjadi inti kajian penelitian.
Persamaan geodesik diturunkan secara eksplisit dari Lagrangian metrik Kerr
melalui metode Euler-Lagrange, dan persamaan transportasi paralel spin
diformulasikan dalam bentuk sistem matriks.
Selanjutnya, metrik Kerr direduksi ke aproksimasi medan lemah
untuk memperoleh ekspresi analitik laju presesi Lense-Thirring.
Sistem persamaan diferensial yang dihasilkan kemudian direduksi ke orde satu
dan diselesaikan secara numerik.


% ============================================================
\section{Desain Penelitian}
\label{sec:desain}

Penelitian ini bersifat kajian analitik-numerik.
Tahapan penelitian meliputi:
\begin{enumerate}
    \item Derivasi persamaan geodesik dari Lagrangian metrik Kerr.
    \item Formulasi persamaan transportasi paralel tensor spin
          dalam bentuk sistem matriks.
    \item Reduksi metrik Kerr ke aproksimasi medan lemah
          untuk menurunkan ekspresi analitik laju presesi Lense-Thirring.
    \item Reduksi sistem persamaan diferensial orde dua ke sistem orde satu.
    \item Penyelesaian numerik menggunakan metode Runge-Kutta orde empat
          yang diimplementasikan dalam Python.
    \item Validasi hasil numerik terhadap data eksperimen \textit{Gravity Probe B}.
\end{enumerate}


% ============================================================
% ============================================================
%  SECTION 3.2: DERIVASI GEODESIK DARI LAGRANGIAN KERR
% ============================================================
% ============================================================

\section{Derivasi Persamaan Geodesik dari Lagrangian Metrik Kerr}
\label{sec:derivasi-geodesik}

Persamaan geodesik dapat diturunkan dari prinsip aksi minimal
tanpa memerlukan pengetahuan eksplisit tentang simbol Christoffel.
Metode ini menggunakan Lagrangian
$\mathcal{L} = \tfrac{1}{2}\,g_{\mu\nu}\,\dot{x}^\mu\,\dot{x}^\nu$
dan persamaan Euler-Lagrange.
Pada bagian ini, metrik Kerr disubstitusikan secara eksplisit ke dalam Lagrangian
dan seluruh persamaan gerak diturunkan langkah demi langkah.


\subsection{Lagrangian Metrik Kerr}
\label{subsec:lagrangian-kerr}

Menggunakan elemen garis metrik Kerr (Persamaan~\ref{eq:kerr-line})
dalam koordinat Boyer-Lindquist $(ct, r, \theta, \phi)$
dan notasi $\dot{x}^\mu \equiv dx^\mu / d\tau$,
Lagrangian partikel bermassa dituliskan sebagai:
\begin{equation}
\boxed{
\begin{aligned}
    \mathcal{L}
    &= \frac{1}{2}\,g_{\mu\nu}\,\dot{x}^\mu\,\dot{x}^\nu \\[6pt]
    &= \frac{1}{2}\left(1 - \frac{r_s r}{\Sigma}\right) c^2\dot{t}^2
     + \frac{r_s\, r\, a\sin^2\!\theta}{\Sigma}\; c\,\dot{t}\,\dot{\phi} \\[4pt]
    &\quad - \frac{1}{2}\,\frac{\Sigma}{\Delta}\,\dot{r}^2
     - \frac{1}{2}\,\Sigma\,\dot{\theta}^2 \\[4pt]
    &\quad - \frac{1}{2}\left(r^2 + a^2
     + \frac{r_s\, r\, a^2\sin^2\!\theta}{\Sigma}\right)\sin^2\!\theta\;\dot{\phi}^2,
\end{aligned}
}
    \label{eq:lagrangian-kerr}
\end{equation}
dengan $\Sigma = r^2 + a^2\cos^2\!\theta$
dan $\Delta = r^2 - r_s r + a^2$ sebagaimana didefinisikan
pada Persamaan~\eqref{eq:sigma-delta}.

Lagrangian ini merupakan fungsi dari koordinat $(r, \theta)$
dan kecepatan umum $(\dot{t}, \dot{r}, \dot{\theta}, \dot{\phi})$,
tetapi \textbf{tidak bergantung secara eksplisit} pada koordinat $t$ dan $\phi$.
Hal ini merupakan konsekuensi dari dua simetri metrik Kerr:
\begin{enumerate}[label=(\roman*)]
    \item \textbf{Stasioneritas} ($\partial_t g_{\mu\nu} = 0$):
          metrik tidak berubah terhadap translasi waktu.
    \item \textbf{Simetri aksial} ($\partial_\phi g_{\mu\nu} = 0$):
          metrik tidak berubah terhadap rotasi di sekitar sumbu simetri.
\end{enumerate}


\subsection{Koordinat Siklik dan Konstanta Gerak}
\label{subsec:konstanta-gerak}

Karena $t$ dan $\phi$ tidak muncul secara eksplisit dalam $\mathcal{L}$,
keduanya merupakan \textit{koordinat siklik}.
Berdasarkan teorema Noether, momentum konjugat
yang bersesuaian adalah konstanta gerak sepanjang geodesik.

\paragraph{Konstanta gerak dari $t$ (energi spesifik).}
Momentum konjugat terhadap $t$ adalah:
\begin{equation}
    p_t \equiv \frac{\partial\mathcal{L}}{\partial(c\dot{t})}
    = g_{tt}\,c\dot{t} + g_{t\phi}\,\dot{\phi}.
    \label{eq:pt-def}
\end{equation}
Substitusi komponen metrik Kerr:
\begin{equation}
    p_t = \left(1 - \frac{r_s r}{\Sigma}\right) c\dot{t}
        + \frac{r_s\, r\, a\sin^2\!\theta}{\Sigma}\;\dot{\phi}.
    \label{eq:pt-explicit}
\end{equation}
Karena $\partial\mathcal{L}/\partial t = 0$,
persamaan Euler-Lagrange $\frac{d}{d\tau}p_t = 0$ menjamin
bahwa $p_t$ konstan sepanjang geodesik.
Konstanta ini diidentifikasi sebagai energi spesifik:
\begin{equation}
    \boxed{p_t \equiv \frac{E}{c} = \text{konstan}.}
    \label{eq:energy-const}
\end{equation}

\paragraph{Konstanta gerak dari $\phi$ (momentum sudut spesifik).}
Momentum konjugat terhadap $\phi$ adalah:
\begin{equation}
    p_\phi \equiv \frac{\partial\mathcal{L}}{\partial\dot{\phi}}
    = g_{\phi t}\,c\dot{t} + g_{\phi\phi}\,\dot{\phi}.
    \label{eq:pphi-def}
\end{equation}
Substitusi komponen metrik:
\begin{equation}
    p_\phi = \frac{r_s\, r\, a\sin^2\!\theta}{\Sigma}\;c\dot{t}
           - \left(r^2 + a^2
             + \frac{r_s\, r\, a^2\sin^2\!\theta}{\Sigma}\right)\sin^2\!\theta\;\dot{\phi}.
    \label{eq:pphi-explicit}
\end{equation}
Karena $\partial\mathcal{L}/\partial\phi = 0$:
\begin{equation}
    \boxed{p_\phi \equiv -L_z = \text{konstan},}
    \label{eq:Lz-const}
\end{equation}
di mana $L_z$ adalah komponen momentum sudut terhadap sumbu rotasi.
Tanda negatif merupakan konvensi yang lazim dalam literatur \parencite{Carroll2004SpacetimeGeometry}.


\subsection{Persamaan Euler-Lagrange untuk Koordinat $r$}
\label{subsec:euler-lagrange-r}

Untuk koordinat $r$, persamaan Euler-Lagrange memberikan:
\begin{equation}
    \frac{d}{d\tau}\frac{\partial\mathcal{L}}{\partial\dot{r}}
    - \frac{\partial\mathcal{L}}{\partial r} = 0.
    \label{eq:el-r}
\end{equation}

Pertama-tama, turunan parsial Lagrangian terhadap kecepatan umum $\dot{r}$ dievaluasi.
Dari Persamaan~\eqref{eq:lagrangian-kerr},
satu-satunya suku yang mengandung $\dot{r}$ adalah
$-\tfrac{1}{2}(\Sigma/\Delta)\dot{r}^2$, sehingga:
\begin{equation}
    \frac{\partial\mathcal{L}}{\partial\dot{r}}
    = -\frac{\Sigma}{\Delta}\,\dot{r}.
    \label{eq:dLdrdot}
\end{equation}

Selanjutnya, turunan terhadap waktu proper dari hasil di atas
dihitung dengan menerapkan aturan produk dan aturan rantai:
\begin{equation}
    \frac{d}{d\tau}\!\left(-\frac{\Sigma}{\Delta}\,\dot{r}\right)
    = -\frac{\Sigma}{\Delta}\,\ddot{r}
      - \frac{d}{d\tau}\!\left(\frac{\Sigma}{\Delta}\right)\dot{r}.
    \label{eq:ddtau-dLdrdot}
\end{equation}
Perhatikan bahwa $\Sigma = r^2 + a^2\cos^2\!\theta$ bergantung pada $r$ dan $\theta$,
sedangkan $\Delta = r^2 - r_s r + a^2$ bergantung pada $r$ saja, sehingga:
\begin{equation}
    \frac{d}{d\tau}\!\left(\frac{\Sigma}{\Delta}\right)
    = \frac{\dot{r}\,\partial_r(\Sigma/\Delta)
          + \dot{\theta}\,\partial_\theta(\Sigma/\Delta)}{1}
    = \frac{\Delta\,(2r\dot{r} - 2a^2\sin\theta\cos\theta\,\dot{\theta})
          - \Sigma\,(2r - r_s)\dot{r}}{\Delta^2}.
    \label{eq:ddt-sigma-delta}
\end{equation}

Di sisi lain, turunan parsial Lagrangian terhadap koordinat $r$ perlu dihitung secara terpisah.
Turunan parsial Lagrangian terhadap $r$ melibatkan seluruh suku
yang mengandung $\Sigma(r,\theta)$ dan $\Delta(r)$.
Turunan parsial yang diperlukan adalah:
\begin{align}
    \frac{\partial\Sigma}{\partial r} &= 2r, \label{eq:dsigma-dr} \\[4pt]
    \frac{\partial\Delta}{\partial r} &= 2r - r_s, \label{eq:ddelta-dr} \\[4pt]
    \frac{\partial}{\partial r}\!\left(\frac{r_s r}{\Sigma}\right)
    &= \frac{r_s(\Sigma - 2r^2)}{\Sigma^2}
     = \frac{r_s(a^2\cos^2\!\theta - r^2)}{\Sigma^2}. \label{eq:drsr-sigma}
\end{align}
Masing-masing suku Lagrangian diturunkan terhadap $r$:
\begin{equation}
\begin{aligned}
    \frac{\partial\mathcal{L}}{\partial r}
    &= \frac{1}{2}\,\frac{\partial g_{tt}}{\partial r}\,c^2\dot{t}^2
     + \frac{\partial g_{t\phi}}{\partial r}\,c\dot{t}\dot{\phi}
     - \frac{1}{2}\,\frac{\partial}{\partial r}\!\left(\frac{\Sigma}{\Delta}\right)\dot{r}^2 \\
    &\quad - \frac{1}{2}\,\frac{\partial\Sigma}{\partial r}\,\dot{\theta}^2
     - \frac{1}{2}\,\frac{\partial g_{\phi\phi}}{\partial r}\,\dot{\phi}^2.
\end{aligned}
    \label{eq:dL-dr}
\end{equation}
Substitusi turunan-turunan di atas memberikan $\partial\mathcal{L}/\partial r$
dalam bentuk eksplisit yang bergantung pada $r$, $\theta$, $\dot{t}$, $\dot{\phi}$, dan $\dot{r}$.

Dengan mensubstitusikan seluruh hasil di atas ke dalam Persamaan~\eqref{eq:el-r}
dan mengalikan kedua ruas dengan $-\Delta/\Sigma$, percepatan radial dapat diisolasi:
\begin{equation}
    \boxed{\ddot{r} = -\Gamma^r_{\mu\nu}\,\dot{x}^\mu\,\dot{x}^\nu
    = f_r(r, \theta, \dot{t}, \dot{r}, \dot{\theta}, \dot{\phi}),}
    \label{eq:geodesic-r}
\end{equation}
di mana $f_r$ merupakan fungsi yang bergantung pada posisi dan kecepatan,
yang secara implisit mengandung seluruh komponen
$\Gamma^r_{\mu\nu}$ dari metrik Kerr.
Bentuk eksplisit $f_r$ sangat panjang dan dalam implementasi praktis
dihitung secara simbolik atau numerik menggunakan pustaka \texttt{EinsteinPy}.


\subsection{Persamaan Euler-Lagrange untuk Koordinat $\theta$}
\label{subsec:euler-lagrange-theta}

Untuk koordinat $\theta$, prosedur yang sama diterapkan:
\begin{equation}
    \frac{d}{d\tau}\frac{\partial\mathcal{L}}{\partial\dot{\theta}}
    - \frac{\partial\mathcal{L}}{\partial\theta} = 0.
    \label{eq:el-theta}
\end{equation}

Turunan parsial terhadap kecepatan umum $\dot{\theta}$ memberikan:
\begin{equation}
    \frac{\partial\mathcal{L}}{\partial\dot{\theta}}
    = -\Sigma\,\dot{\theta}.
    \label{eq:dLdthetadot}
\end{equation}

Turunan terhadap waktu proper dari ekspresi tersebut menghasilkan:
\begin{equation}
    \frac{d}{d\tau}\!\left(-\Sigma\,\dot{\theta}\right)
    = -\Sigma\,\ddot{\theta}
      - (2r\dot{r} - 2a^2\sin\theta\cos\theta\,\dot{\theta})\,\dot{\theta}.
    \label{eq:ddtau-dLdthetadot}
\end{equation}

Sementara itu, turunan parsial Lagrangian terhadap $\theta$ memerlukan
perhitungan yang lebih teliti.
Turunan parsial ini lebih kompleks karena banyak komponen metrik
yang bergantung pada $\theta$ melalui $\Sigma$ dan $\sin^2\!\theta$.
Turunan yang relevan meliputi:
\begin{align}
    \frac{\partial\Sigma}{\partial\theta}
    &= -2a^2\sin\theta\cos\theta, \label{eq:dsigma-dtheta} \\[4pt]
    \frac{\partial}{\partial\theta}\!\left(\frac{r_s r}{\Sigma}\right)
    &= \frac{2r_s\, r\, a^2\sin\theta\cos\theta}{\Sigma^2}, \label{eq:drsr-dtheta} \\[4pt]
    \frac{\partial(\sin^2\!\theta)}{\partial\theta}
    &= 2\sin\theta\cos\theta. \label{eq:dsin2-dtheta}
\end{align}
Dengan demikian:
\begin{equation}
\begin{aligned}
    \frac{\partial\mathcal{L}}{\partial\theta}
    &= \frac{1}{2}\,\frac{\partial g_{tt}}{\partial\theta}\,c^2\dot{t}^2
     + \frac{\partial g_{t\phi}}{\partial\theta}\,c\dot{t}\dot{\phi}
     - \frac{1}{2}\,\frac{\partial}{\partial\theta}\!\left(\frac{\Sigma}{\Delta}\right)\dot{r}^2 \\
    &\quad + a^2\sin\theta\cos\theta\,\dot{\theta}^2
     - \frac{1}{2}\,\frac{\partial g_{\phi\phi}}{\partial\theta}\,\dot{\phi}^2.
\end{aligned}
    \label{eq:dL-dtheta}
\end{equation}

Penggabungan seluruh suku dan isolasi $\ddot{\theta}$
menghasilkan persamaan gerak polar:
\begin{equation}
    \boxed{\ddot{\theta} = -\Gamma^\theta_{\mu\nu}\,\dot{x}^\mu\,\dot{x}^\nu
    = f_\theta(r, \theta, \dot{t}, \dot{r}, \dot{\theta}, \dot{\phi}).}
    \label{eq:geodesic-theta}
\end{equation}


\subsection{Penyelesaian $\dot{t}$ dan $\dot{\phi}$ dari Konstanta Gerak}
\label{subsec:solve-tdot-phidot}

Persamaan~\eqref{eq:pt-explicit} dan \eqref{eq:pphi-explicit}
membentuk sistem linier dua persamaan dengan dua variabel,
$c\dot{t}$ dan $\dot{\phi}$.
Sistem ini dapat dituliskan dalam bentuk matriks:
\begin{equation}
    \begin{pmatrix} g_{tt} & g_{t\phi} \\ g_{\phi t} & g_{\phi\phi} \end{pmatrix}
    \begin{pmatrix} c\dot{t} \\ \dot{\phi} \end{pmatrix}
    =
    \begin{pmatrix} E/c \\ -L_z \end{pmatrix}.
    \label{eq:system-tdot-phidot}
\end{equation}
Solusinya diperoleh melalui inversi matriks $2 \times 2$:
\begin{equation}
    \begin{pmatrix} c\dot{t} \\ \dot{\phi} \end{pmatrix}
    = \frac{1}{g_{tt}g_{\phi\phi} - g_{t\phi}^2}
    \begin{pmatrix} g_{\phi\phi} & -g_{t\phi} \\ -g_{\phi t} & g_{tt} \end{pmatrix}
    \begin{pmatrix} E/c \\ -L_z \end{pmatrix}.
    \label{eq:solution-tdot-phidot}
\end{equation}
Determinan sub-matriks $t$-$\phi$ dapat dihitung dari komponen metrik Kerr:
\begin{equation}
    g_{tt}g_{\phi\phi} - g_{t\phi}^2 = -\Delta\sin^2\!\theta,
    \label{eq:det-tphi}
\end{equation}
sehingga:
\begin{align}
    c\dot{t} &= \frac{1}{\Delta\sin^2\!\theta}
    \left[
        \left(r^2 + a^2 + \frac{r_s\, r\, a^2\sin^2\!\theta}{\Sigma}\right)
        \sin^2\!\theta\;\frac{E}{c}
        - \frac{r_s\, r\, a\sin^2\!\theta}{\Sigma}\; L_z
    \right],
    \label{eq:tdot-explicit} \\[8pt]
    \dot{\phi} &= \frac{1}{\Delta\sin^2\!\theta}
    \left[
        \frac{r_s\, r\, a\sin^2\!\theta}{\Sigma}\;\frac{E}{c}
        + \left(1 - \frac{r_s r}{\Sigma}\right) L_z
    \right].
    \label{eq:phidot-explicit}
\end{align}

Dengan demikian, evolusi temporal $t(\tau)$ dan $\phi(\tau)$
ditentukan sepenuhnya oleh konstanta gerak $E$ dan $L_z$
serta posisi sesaat $(r, \theta)$.
Persamaan \eqref{eq:geodesic-r} dan \eqref{eq:geodesic-theta}
bersama dengan \eqref{eq:tdot-explicit}--\eqref{eq:phidot-explicit}
membentuk sistem lengkap persamaan gerak pada metrik Kerr.


\subsection{Constraint Normalisasi}
\label{subsec:constraint}

Sebagai verifikasi konsistensi, normalisasi empat-kecepatan harus dipenuhi
sepanjang lintasan geodesik:
\begin{equation}
    g_{\mu\nu}\,\dot{x}^\mu\,\dot{x}^\nu = c^2.
    \label{eq:norm-constraint}
\end{equation}
Substitusi metrik Kerr memberikan:
\begin{equation}
    \left(1 - \frac{r_s r}{\Sigma}\right) c^2\dot{t}^2
    + \frac{2\,r_s\, r\, a\sin^2\!\theta}{\Sigma}\, c\dot{t}\dot{\phi}
    - \frac{\Sigma}{\Delta}\,\dot{r}^2
    - \Sigma\,\dot{\theta}^2
    - \left(r^2 + a^2 + \frac{r_s\, r\, a^2\sin^2\!\theta}{\Sigma}\right)
      \sin^2\!\theta\,\dot{\phi}^2 = c^2.
    \label{eq:norm-explicit}
\end{equation}
Persamaan ini tidak independen dari persamaan geodesik ---
ia merupakan integral pertama gerak.
Dalam implementasi numerik, deviasi dari nilai $c^2$ digunakan
sebagai indikator kestabilan dan akurasi integrasi
(lihat Bab~\ref{chap:hasil}).


% ============================================================
% ============================================================
%  SECTION 3.3: TRANSPORTASI PARALEL SPIN — BENTUK MATRIKS
% ============================================================
% ============================================================

\section{Transportasi Paralel Spin dalam Bentuk Matriks}
\label{sec:spin-matriks}

Persamaan transportasi paralel spin (Persamaan~\ref{eq:spin-transport})
mendeskripsikan evolusi vektor spin $S^\mu$ sepanjang geodesik.
Pada bagian ini, persamaan tersebut diformulasikan secara eksplisit
sebagai sistem ODE orde satu dalam bentuk matriks,
yang sesuai untuk penyelesaian numerik.


\subsection{Formulasi Sistem ODE}
\label{subsec:spin-ode-system}

Dengan mendefinisikan empat-kecepatan $u^\rho(\tau) = \dot{x}^\rho(\tau)$
yang diperoleh dari penyelesaian geodesik,
Persamaan~\eqref{eq:spin-ode} dapat ditulis komponen per komponen:
\begin{equation}
\begin{aligned}
    \frac{dS^t}{d\tau}
    &= -\left(\Gamma^t_{tt}\, u^t + \Gamma^t_{tr}\, u^r
      + \Gamma^t_{t\theta}\, u^\theta + \Gamma^t_{t\phi}\, u^\phi
      + \Gamma^t_{r\phi}\, u^\phi + \cdots\right) S^t \\
    &\quad -\left(\Gamma^t_{rt}\, u^t + \Gamma^t_{rr}\, u^r + \cdots\right) S^r
     - \cdots \\[4pt]
    &= -\sum_{\nu=0}^{3} A^t_{\ \nu}(\tau)\, S^\nu,
\end{aligned}
    \label{eq:spin-t-component}
\end{equation}
dengan $A^\mu_{\ \nu}(\tau) = \sum_\rho \Gamma^\mu_{\nu\rho}\, u^\rho(\tau)$.

Secara umum, sistem lengkap ditulis sebagai:
\begin{equation}
    \frac{d}{d\tau}
    \begin{pmatrix} S^t \\ S^r \\ S^\theta \\ S^\phi \end{pmatrix}
    = -\underbrace{
    \begin{pmatrix}
        A^t_{\ t} & A^t_{\ r} & A^t_{\ \theta} & A^t_{\ \phi} \\
        A^r_{\ t} & A^r_{\ r} & A^r_{\ \theta} & A^r_{\ \phi} \\
        A^\theta_{\ t} & A^\theta_{\ r} & A^\theta_{\ \theta} & A^\theta_{\ \phi} \\
        A^\phi_{\ t} & A^\phi_{\ r} & A^\phi_{\ \theta} & A^\phi_{\ \phi}
    \end{pmatrix}
    }_{\displaystyle\mathbf{A}(\tau)}
    \begin{pmatrix} S^t \\ S^r \\ S^\theta \\ S^\phi \end{pmatrix}.
    \label{eq:spin-full-matrix}
\end{equation}
Matriks $\mathbf{A}(\tau)$ berukuran $4 \times 4$
dan bergantung pada waktu melalui posisi $x^\mu(\tau)$
dan kecepatan $u^\mu(\tau)$ yang merupakan solusi geodesik.
Elemen-elemen matriks ini dihitung dari:
\begin{equation}
    \boxed{A^\mu_{\ \nu}(\tau) = \sum_{\rho=0}^{3} \Gamma^\mu_{\nu\rho}\bigl(x(\tau)\bigr)\; u^\rho(\tau).}
    \label{eq:A-elements}
\end{equation}


\subsection{Contoh Elemen Matriks}
\label{subsec:contoh-elemen-A}

Sebagai ilustrasi, beberapa elemen matriks $\mathbf{A}$ dapat ditulis secara eksplisit:
\begin{align}
    A^t_{\ t} &= \Gamma^t_{tt}\,u^t + \Gamma^t_{tr}\,u^r
              + \Gamma^t_{t\theta}\,u^\theta + \Gamma^t_{t\phi}\,u^\phi,
    \label{eq:Att} \\[4pt]
    A^t_{\ \phi} &= \Gamma^t_{\phi t}\,u^t + \Gamma^t_{\phi r}\,u^r
                 + \Gamma^t_{\phi\theta}\,u^\theta + \Gamma^t_{\phi\phi}\,u^\phi,
    \label{eq:Atphi} \\[4pt]
    A^r_{\ r} &= \Gamma^r_{rt}\,u^t + \Gamma^r_{rr}\,u^r
              + \Gamma^r_{r\theta}\,u^\theta + \Gamma^r_{r\phi}\,u^\phi,
    \label{eq:Arr} \\[4pt]
    A^\phi_{\ t} &= \Gamma^\phi_{tt}\,u^t + \Gamma^\phi_{tr}\,u^r
                 + \Gamma^\phi_{t\theta}\,u^\theta + \Gamma^\phi_{t\phi}\,u^\phi.
    \label{eq:Aphit}
\end{align}
Seluruh simbol Christoffel $\Gamma^\mu_{\nu\rho}$ dievaluasi pada posisi sesaat $x^\mu(\tau)$
dan dihitung dari definisi (Persamaan~\ref{eq:christoffel-def})
dengan metrik Kerr.

Dalam implementasi numerik, matriks $\mathbf{A}(\tau)$ dievaluasi ulang
pada setiap langkah integrasi bersamaan dengan penyelesaian geodesik.
Perhitungan ini dilakukan secara otomatis oleh pustaka \texttt{EinsteinPy}
yang menyediakan fungsi untuk menghitung seluruh komponen Christoffel
dari metrik yang diberikan \parencite{Bapat2020EinsteinPy}.


\subsection{Constraint Ortogonalitas}
\label{subsec:constraint-ortogonalitas}

Vektor spin harus memenuhi constraint ortogonalitas terhadap empat-kecepatan
pada setiap titik sepanjang geodesik:
\begin{equation}
    g_{\mu\nu}\, S^\mu\, u^\nu = 0.
    \label{eq:ortogonalitas}
\end{equation}
Dalam bentuk eksplisit untuk metrik Kerr:
\begin{equation}
    g_{tt}\, S^t\, u^t + g_{t\phi}(S^t\, u^\phi + S^\phi\, u^t)
    + g_{rr}\, S^r\, u^r + g_{\theta\theta}\, S^\theta\, u^\theta
    + g_{\phi\phi}\, S^\phi\, u^\phi = 0.
    \label{eq:ortogonalitas-eksplisit}
\end{equation}

Constraint ini bersifat integral pertama:
jika dipenuhi pada kondisi awal $\tau = 0$,
maka transportasi paralel secara otomatis melestarikannya
untuk seluruh $\tau > 0$.
Dalam praktik numerik, deviasi dari nol pada kuantitas
$g_{\mu\nu} S^\mu u^\nu$ digunakan sebagai indikator
akurasi integrasi, serupa dengan constraint normalisasi pada geodesik
(Persamaan~\ref{eq:norm-constraint}).

Constraint ini juga memiliki interpretasi fisik yang penting:
ia menjamin bahwa vektor spin bersifat ``ruang murni'' (\textit{spacelike})
dalam kerangka diam sesaat partikel.
Artinya, komponen spin tidak memiliki proyeksi terhadap arah waktu
dalam kerangka acuan ko-bergerak (\textit{comoving frame}) giroskop
\parencite{Misner1973Gravitation}.


% ============================================================
% ============================================================
%  SECTION 3.4: REDUKSI METRIK KERR KE APROKSIMASI MEDAN LEMAH
% ============================================================
% ============================================================

\section{Reduksi Metrik Kerr ke Aproksimasi Medan Lemah}
\label{sec:medan-lemah}

Pada bagian ini ditunjukkan secara eksplisit bagaimana metrik Kerr
mereduksi menjadi bentuk linier pada limit medan lemah dan rotasi lambat,
serta bagaimana ekspresi analitik laju presesi Lense-Thirring diperoleh
dari persamaan transportasi paralel pada limit tersebut.


\subsection{Syarat Aproksimasi}
\label{subsec:syarat-aproksimasi}

Untuk medan gravitasi Bumi, dua parameter tanpa dimensi berikut bernilai sangat kecil:
\begin{equation}
    \epsilon_1 \equiv \frac{r_s}{r} = \frac{2GM}{c^2 r}
    \approx \frac{8{,}87 \times 10^{-3}}{7{,}013 \times 10^{6}}
    \approx 1{,}26 \times 10^{-9} \ll 1,
    \label{eq:eps1}
\end{equation}
\begin{equation}
    \epsilon_2 \equiv \frac{a}{r} = \frac{J}{Mcr}
    \approx \frac{3{,}26}{7{,}013 \times 10^{6}}
    \approx 4{,}65 \times 10^{-7} \ll 1.
    \label{eq:eps2}
\end{equation}
Aproksimasi medan lemah (\textit{weak-field limit})
mensyaratkan bahwa seluruh suku orde $\mathcal{O}(\epsilon_1^2)$,
$\mathcal{O}(\epsilon_2^2)$, dan $\mathcal{O}(\epsilon_1\epsilon_2)$
dapat diabaikan terhadap suku orde pertama.


\subsection{Ekspansi Komponen Metrik Kerr}
\label{subsec:ekspansi-metrik}

Setiap komponen metrik Kerr (Persamaan~\ref{eq:kerr-matrix})
diekspansi dalam deret Taylor terhadap $\epsilon_1$ dan $\epsilon_2$.

Ekspansi dimulai dari fungsi $\Sigma^{-1}$.
Karena $\Sigma = r^2 + a^2\cos^2\!\theta = r^2(1 + \epsilon_2^2\cos^2\!\theta)$, diperoleh:
\begin{equation}
    \frac{1}{\Sigma} = \frac{1}{r^2}\,\frac{1}{1 + \epsilon_2^2\cos^2\!\theta}
    = \frac{1}{r^2}\left(1 - \epsilon_2^2\cos^2\!\theta + \mathcal{O}(\epsilon_2^4)\right)
    \approx \frac{1}{r^2}.
    \label{eq:sigma-inv-expand}
\end{equation}

Dengan cara serupa, fungsi $\Delta$ diekspansi.
Karena $\Delta = r^2 - r_s r + a^2 = r^2(1 - \epsilon_1 + \epsilon_2^2 r^{-2} a^2)$:
\begin{equation}
    \Delta \approx r^2(1 - \epsilon_1) = r^2 - r_s r.
    \label{eq:delta-expand}
\end{equation}

Komponen temporal $g_{tt}$ pada limit ini menjadi:
\begin{equation}
    g_{tt} = 1 - \frac{r_s r}{\Sigma}
    \approx 1 - \frac{r_s r}{r^2}
    = 1 - \frac{r_s}{r}
    = 1 - \frac{2GM}{c^2 r}.
    \label{eq:gtt-wf}
\end{equation}

Untuk komponen radial $g_{rr}$, ekspansi memberikan:
\begin{equation}
    g_{rr} = -\frac{\Sigma}{\Delta}
    \approx -\frac{r^2}{r^2 - r_s r}
    = -\frac{1}{1 - r_s/r}
    \approx -\left(1 + \frac{r_s}{r}\right)
    = -\left(1 + \frac{2GM}{c^2 r}\right).
    \label{eq:grr-wf}
\end{equation}

Komponen diagonal spasial $g_{\theta\theta}$ dan $g_{\phi\phi}$
tereduksi menjadi bentuk ruang datar:
\begin{align}
    g_{\theta\theta} &= -\Sigma \approx -r^2,
    \label{eq:gthetatheta-wf} \\[4pt]
    g_{\phi\phi}^{(\text{diag})} &= -\left(r^2 + a^2
    + \frac{r_s r a^2\sin^2\!\theta}{\Sigma}\right)\sin^2\!\theta
    \approx -r^2\sin^2\!\theta.
    \label{eq:gphiphi-wf}
\end{align}

Komponen yang paling signifikan secara fisik adalah suku campuran $g_{t\phi}$,
yang merupakan satu-satunya pembawa efek \textit{frame-dragging}:
\begin{equation}
    g_{t\phi} = \frac{r_s\, r\, a\sin^2\!\theta}{\Sigma}
    \approx \frac{r_s\, a\sin^2\!\theta}{r}
    = \frac{2GM\,a\sin^2\!\theta}{c^2\, r}
    = \frac{2GJ\sin^2\!\theta}{c^3\, r}.
    \label{eq:gtphi-wf}
\end{equation}
Di sini digunakan hubungan $a = J/(Mc)$, sehingga $r_s\, a = 2GM\,a/c^2 = 2GJ/c^3$.

Komponen ini merupakan \textbf{satu-satunya suku yang bergantung
pada momentum sudut $J$} dan \textbf{lenyap jika $J = 0$} (massa tak berotasi).
Suku inilah yang secara langsung bertanggung jawab atas efek Lense-Thirring.


\subsection{Metrik Linierisasi (Weak-Field Kerr)}
\label{subsec:metrik-linearisasi}

Menggabungkan seluruh hasil ekspansi, metrik Kerr pada limit medan lemah
dapat dituliskan dalam bentuk matriks $4 \times 4$ yang terlinearisasi:
\begin{equation}
    g_{\mu\nu}^{(\text{wf})} \approx
    \begin{pmatrix}
        1 - \dfrac{2GM}{c^2 r} & 0 & 0 & \dfrac{2GJ\sin^2\!\theta}{c^3 r} \\[10pt]
        0 & -\!\left(1 + \dfrac{2GM}{c^2 r}\right) & 0 & 0 \\[10pt]
        0 & 0 & -r^2 & 0 \\[10pt]
        \dfrac{2GJ\sin^2\!\theta}{c^3 r} & 0 & 0 & -r^2\sin^2\!\theta
    \end{pmatrix}.
    \label{eq:kerr-wf-matrix}
\end{equation}
Metrik ini identik dengan metrik Lense-Thirring yang secara historis
diturunkan langsung dari persamaan medan Einstein terlinearisasi
\parencite{Lense1918LenseThirring}.
Perbedaan terhadap metrik Schwarzschild hanya terletak pada
komponen campuran $g_{t\phi} = g_{\phi t} = 2GJ\sin^2\!\theta/(c^3 r)$.


\subsection{Derivasi Laju Presesi Lense-Thirring dari Transportasi Spin}
\label{subsec:derivasi-lt}

Pada bagian ini ditunjukkan bagaimana persamaan transportasi paralel
spin (Persamaan~\ref{eq:spin-full-matrix}) mereduksi menjadi
persamaan presesi sederhana pada limit medan lemah,
yang menghasilkan ekspresi laju Lense-Thirring.

Pada limit medan lemah, partikel uji bergerak dengan kecepatan
jauh di bawah kecepatan cahaya ($v \ll c$),
sehingga komponen empat-kecepatan didominasi oleh:
\begin{equation}
    u^\mu \approx (c,\; v^r,\; v^\theta/r,\; v^\phi/(r\sin\theta)),
    \label{eq:u-nonrel}
\end{equation}
di mana $v^i$ adalah komponen kecepatan tiga-dimensi,
dan $u^t \approx c$ karena faktor Lorentz $\gamma \approx 1$.

Dari metrik terlinearisasi (Persamaan~\ref{eq:kerr-wf-matrix}),
simbol Christoffel kemudian dihitung menggunakan Persamaan~\eqref{eq:christoffel-def}.
Kontribusi dari suku $g_{t\phi}$ terhadap Christoffel yang relevan
untuk transportasi spin adalah:
\begin{align}
    \Gamma^i_{t j} &= \frac{1}{2}\,g^{ik}
    \left(\partial_j g_{kt} + \partial_t g_{kj} - \partial_k g_{tj}\right)
    \approx \frac{1}{2}\,\delta^{ik}
    \left(\partial_j g_{kt} - \partial_k g_{tj}\right),
    \label{eq:christoffel-wf-spatial}
\end{align}
di mana indeks Latin $i,j,k \in \{r, \theta, \phi\}$
menandai komponen spasial, dan $g^{ik} \approx -\delta^{ik}/h_{ii}$
dengan $h_{ii}$ adalah komponen metrik spasial diagonal.
Suku $\partial_t g_{kj} = 0$ karena metrik stasioner.

Untuk menyederhanakan analisis, didefinisikan potensial gravitomagnetik $\mathbf{A}_g$
yang diekstraksi dari komponen campuran metrik:
\begin{equation}
    (A_g)_i \equiv \frac{c}{2}\, g_{ti}.
    \label{eq:gravitomagnetic-potential}
\end{equation}
Untuk metrik Kerr medan lemah, satu-satunya komponen non-nol adalah:
\begin{equation}
    (A_g)_\phi = \frac{c}{2}\, g_{t\phi}
    = \frac{GJ\sin^2\!\theta}{c^2 r}.
    \label{eq:Ag-phi}
\end{equation}
Dalam koordinat Kartesian $(x, y, z)$, potensial ini ekuivalen dengan:
\begin{equation}
    \mathbf{A}_g = \frac{G}{c^2}\,\frac{\mathbf{J} \times \mathbf{r}}{r^3}.
    \label{eq:Ag-cartesian}
\end{equation}

Secara analog dengan elektromagnetisme, medan gravitomagnetik
$\mathbf{B}_g$ diperoleh sebagai rotasi dari potensial tersebut:
\begin{equation}
    \mathbf{B}_g \equiv \nabla \times \mathbf{A}_g.
    \label{eq:Bg-def}
\end{equation}
Substitusi Persamaan~\eqref{eq:Ag-cartesian} dan evaluasi rotasi
menggunakan identitas vektor
$\nabla \times \left(\frac{\mathbf{J} \times \mathbf{r}}{r^3}\right)
= \frac{1}{r^3}\left[3\hat{\mathbf{r}}(\hat{\mathbf{r}} \cdot \mathbf{J})
- \mathbf{J}\right] - \frac{8\pi}{3}\,\mathbf{J}\,\delta^3(\mathbf{r})$,
untuk $r > 0$ diperoleh:
\begin{equation}
    \mathbf{B}_g = \frac{G}{c^2 r^3}
    \left[3\hat{\mathbf{r}}\,(\hat{\mathbf{r}} \cdot \mathbf{J}) - \mathbf{J}\right].
    \label{eq:Bg-explicit}
\end{equation}

Dengan medan gravitomagnetik yang telah diperoleh,
persamaan transportasi paralel spin
$dS^i/d\tau = -\Gamma^i_{\nu\rho}\, S^\nu\, u^\rho$
(Persamaan~\ref{eq:spin-ode}) untuk komponen spasial
kini tereduksi menjadi persamaan presesi yang familiar:
\begin{equation}
    \frac{d\mathbf{S}}{dt} = \bm{\Omega} \times \mathbf{S},
    \label{eq:precession-eq}
\end{equation}
di mana vektor kecepatan sudut presesi $\bm{\Omega}$ merupakan
superposisi kontribusi geodesik dan Lense-Thirring.
Kontribusi Lense-Thirring diperoleh dari kopling antara spin
dan medan gravitomagnetik \parencite{Will2014Confrontation}:
\begin{equation}
    \bm{\Omega}_{\text{LT}} = \frac{1}{2}\,\mathbf{B}_g
    = \frac{G}{2c^2 r^3}
    \left[3\hat{\mathbf{r}}\,(\hat{\mathbf{r}} \cdot \mathbf{J}) - \mathbf{J}\right].
    \label{eq:omega-lt-from-Bg}
\end{equation}
Ekspresi ini identik dengan Persamaan~\eqref{eq:omega-lt-general}
yang diperoleh dari pendekatan langsung pada Bab~\ref{chap:tinjauan}.

Khusus untuk orbit polar (inklinasi $\sim 90°$)
seperti yang ditempuh satelit GP-B,
vektor posisi rata-rata hampir tegak lurus terhadap sumbu rotasi Bumi,
sehingga $\hat{\mathbf{r}} \cdot \mathbf{J} \approx 0$ secara rata-rata orbit.
Pada kondisi ini, Persamaan~\eqref{eq:omega-lt-from-Bg} mereduksi menjadi:
\begin{equation}
    \bm{\Omega}_{\text{LT}} \approx -\frac{G\,\mathbf{J}}{2c^2 r^3},
    \label{eq:omega-lt-polar-exact}
\end{equation}
dengan magnitudo:
\begin{equation}
    \boxed{|\bm{\Omega}_{\text{LT}}| \approx \frac{G\,|\mathbf{J}|}{2c^2 r^3}.}
    \label{eq:omega-lt-magnitude-half}
\end{equation}

Persamaan ini berbeda faktor 4 dari aproksimasi yang sering dijumpai
dalam literatur, $\Omega_{\text{LT}} \approx 2GJ/(c^2 r^3)$
(Persamaan~\ref{eq:omega-lt-approx}).
Perbedaan ini berasal dari cara rata-rata orbit dilakukan:
ekspresi $2GJ/(c^2 r^3)$ mencakup kontribusi suku
$3\hat{\mathbf{r}}(\hat{\mathbf{r}} \cdot \mathbf{J})$
yang tidak lenyap secara lokal melainkan hanya setelah rata-rata sepanjang orbit penuh.
Perhitungan yang lebih teliti oleh \textcite{Adler2014}
menunjukkan bahwa laju presesi \textit{frame-dragging}
yang terukur oleh GP-B diberikan oleh rata-rata orbit:
\begin{equation}
    \boxed{
    \langle\bm{\Omega}_{\text{LT}}\rangle_{\text{orbit}}
    = \frac{G}{2c^2 r^3}
    \left[3\hat{\mathbf{L}}\,(\hat{\mathbf{L}} \cdot \mathbf{J}) - \mathbf{J}\right],
    }
    \label{eq:omega-lt-orbit-avg}
\end{equation}
di mana $\hat{\mathbf{L}}$ adalah vektor satuan normal bidang orbit.
Untuk orbit polar GP-B ($\hat{\mathbf{L}} \perp \mathbf{J}$):
\begin{equation}
    |\langle\bm{\Omega}_{\text{LT}}\rangle| = \frac{GJ}{2c^2 r^3}.
    \label{eq:omega-lt-gpb-value}
\end{equation}

Sebagai verifikasi, dilakukan substitusi numerik
menggunakan parameter Bumi dari Tabel~\ref{tab:param-gpb}:
\begin{equation}
\begin{aligned}
    |\Omega_{\text{LT}}|
    &= \frac{G\,J}{2c^2\,r^3} \\[4pt]
    &= \frac{(6{,}674 \times 10^{-11})(5{,}86 \times 10^{33})}
           {2\,(2{,}998 \times 10^8)^2\,(7{,}013 \times 10^6)^3} \\[4pt]
    &= \frac{3{,}91 \times 10^{23}}{2 \times 8{,}99 \times 10^{16} \times 3{,}45 \times 10^{20}} \\[4pt]
    &= \frac{3{,}91 \times 10^{23}}{6{,}20 \times 10^{37}} \\[4pt]
    &\approx 6{,}3 \times 10^{-15}~\text{rad/s}.
\end{aligned}
    \label{eq:omega-lt-numeric}
\end{equation}
Konversi ke mas/yr menggunakan faktor dari Persamaan~\eqref{eq:konversi-masyr}:
\begin{equation}
    |\Omega_{\text{LT}}| \approx 6{,}3 \times 10^{-15} \times 6{,}51 \times 10^{15}
    \approx 41~\text{mas/yr}.
    \label{eq:omega-lt-masyr}
\end{equation}
Nilai ini berada dalam orde yang sama dengan hasil pengukuran GP-B
($37{,}2 \pm 7{,}2$~mas/yr), dan perbedaan minor berasal dari
penyederhanaan orbit sirkular ideal serta nilai parameter Bumi
yang digunakan.
Analisis numerik yang lebih presisi disajikan pada Bab~\ref{chap:hasil}.


\subsection{Derivasi Presesi Geodesik pada Limit Medan Lemah}
\label{subsec:derivasi-geodesik-wf}

Untuk kelengkapan, presesi geodesik (de Sitter) juga diturunkan
pada limit yang sama.
Kontribusi geodesik berasal dari komponen diagonal metrik
(tanpa suku $g_{t\phi}$) dan diberikan oleh \parencite{Will2014Confrontation}:
\begin{equation}
    \bm{\Omega}_{\text{geodetik}}
    = \frac{3}{2}\,\frac{GM}{c^2 r^3}\;\mathbf{r} \times \mathbf{v}.
    \label{eq:omega-geo-wf}
\end{equation}
Untuk orbit sirkular GP-B ($v = \sqrt{GM/r}$, $|\mathbf{r} \times \mathbf{v}| = rv$):
\begin{equation}
\begin{aligned}
    |\Omega_{\text{geodetik}}|
    &= \frac{3}{2}\,\frac{(GM)^{3/2}}{c^2\, r^{5/2}} \\[4pt]
    &= \frac{3}{2}\,\frac{(6{,}674 \times 10^{-11} \times 5{,}972 \times 10^{24})^{3/2}}
           {(2{,}998 \times 10^8)^2\,(7{,}013 \times 10^6)^{5/2}} \\[4pt]
    &\approx 1{,}02 \times 10^{-12}~\text{rad/s} \\[4pt]
    &\approx 6600~\text{mas/yr},
\end{aligned}
    \label{eq:omega-geo-numeric}
\end{equation}
konsisten dengan hasil GP-B ($6602 \pm 18$~mas/yr).


% ============================================================
% ============================================================
%  SECTION 3.5: REDUKSI KE SISTEM ODE ORDE SATU
% ============================================================
% ============================================================

\section{Reduksi ke Sistem Persamaan Diferensial Orde Satu}
\label{sec:reduksi-ode}

Persamaan geodesik (Persamaan~\ref{eq:geodesic}) merupakan
sistem persamaan diferensial biasa (ODE) orde dua:
\begin{equation}
    \ddot{x}^\mu = -\Gamma^\mu_{\nu\rho}\,\dot{x}^\nu\,\dot{x}^\rho,
    \qquad \mu = 0, 1, 2, 3.
    \label{eq:geodesic-ode2}
\end{equation}
Metode integrasi numerik standar (termasuk Runge-Kutta)
dirancang untuk memproses sistem ODE orde satu dalam bentuk
$\dot{\mathbf{Y}} = \mathbf{F}(\tau, \mathbf{Y})$.
Oleh karena itu, sistem orde dua di atas harus direduksi
menjadi sistem orde satu yang ekuivalen.


\subsection{Definisi Variabel Bantu}
\label{subsec:variabel-bantu}

Didefinisikan empat-kecepatan sebagai variabel bantu:
\begin{equation}
    u^\mu \equiv \dot{x}^\mu = \frac{dx^\mu}{d\tau},
    \qquad \mu = 0, 1, 2, 3.
    \label{eq:u-def}
\end{equation}
Dengan substitusi ini, persamaan geodesik orde dua menjadi
dua himpunan persamaan orde satu:
\begin{align}
    \frac{dx^\mu}{d\tau} &= u^\mu,
    \label{eq:ode1-x} \\[6pt]
    \frac{du^\mu}{d\tau} &= -\Gamma^\mu_{\nu\rho}(x)\; u^\nu\, u^\rho.
    \label{eq:ode1-u}
\end{align}
Persamaan~\eqref{eq:ode1-x} berjumlah 4 (evolusi posisi)
dan Persamaan~\eqref{eq:ode1-u} berjumlah 4 (evolusi kecepatan),
menghasilkan total 8 ODE orde satu.


\subsection{Pembentukan Vektor Keadaan}
\label{subsec:state-vector}

Seluruh variabel digabungkan ke dalam vektor keadaan
berdimensi 8:
\begin{equation}
    \boxed{
    \mathbf{Y} =
    \begin{pmatrix}
        x^0 \\ x^1 \\ x^2 \\ x^3 \\ u^0 \\ u^1 \\ u^2 \\ u^3
    \end{pmatrix}
    =
    \begin{pmatrix}
        ct \\ r \\ \theta \\ \phi \\ u^t \\ u^r \\ u^\theta \\ u^\phi
    \end{pmatrix}.
    }
    \label{eq:state-vector}
\end{equation}
Sistem ODE orde satu dituliskan secara kompak sebagai:
\begin{equation}
    \frac{d\mathbf{Y}}{d\tau} = \mathbf{F}(\mathbf{Y}),
    \label{eq:ode-compact}
\end{equation}
dengan fungsi sisi kanan:
\begin{equation}
    \mathbf{F}(\mathbf{Y}) =
    \begin{pmatrix}
        u^t \\[2pt]
        u^r \\[2pt]
        u^\theta \\[2pt]
        u^\phi \\[2pt]
        -\Gamma^t_{\nu\rho}\, u^\nu u^\rho \\[2pt]
        -\Gamma^r_{\nu\rho}\, u^\nu u^\rho \\[2pt]
        -\Gamma^\theta_{\nu\rho}\, u^\nu u^\rho \\[2pt]
        -\Gamma^\phi_{\nu\rho}\, u^\nu u^\rho
    \end{pmatrix}.
    \label{eq:F-explicit}
\end{equation}
Persamaan transportasi spin (Persamaan~\ref{eq:spin-full-matrix})
ditambahkan ke dalam sistem sebagai 4 ODE tambahan:
\begin{equation}
    \frac{dS^\mu}{d\tau} = -A^\mu_{\ \nu}\, S^\nu,
    \qquad A^\mu_{\ \nu} = \Gamma^\mu_{\nu\rho}\, u^\rho,
    \label{eq:spin-ode-recap}
\end{equation}
sehingga vektor keadaan total yang diintegrasikan berdimensi $8 + 4 = 12$:
\begin{equation}
    \mathbf{Y}_{\text{total}} =
    \begin{pmatrix}
        ct \\ r \\ \theta \\ \phi \\
        u^t \\ u^r \\ u^\theta \\ u^\phi \\
        S^t \\ S^r \\ S^\theta \\ S^\phi
    \end{pmatrix}.
    \label{eq:state-vector-total}
\end{equation}


\subsection{Bentuk Eksplisit untuk Setiap Komponen}
\label{subsec:ode-eksplisit}

Secara eksplisit, 12 persamaan ODE orde satu tersebut dituliskan sebagai:
\begin{align}
    \frac{d(ct)}{d\tau} &= u^t, \label{eq:ode-ct}\\[3pt]
    \frac{dr}{d\tau} &= u^r, \label{eq:ode-r}\\[3pt]
    \frac{d\theta}{d\tau} &= u^\theta, \label{eq:ode-theta}\\[3pt]
    \frac{d\phi}{d\tau} &= u^\phi, \label{eq:ode-phi}\\[3pt]
    \frac{du^t}{d\tau} &= -\Gamma^t_{\nu\rho}\, u^\nu\, u^\rho, \label{eq:ode-ut}\\[3pt]
    \frac{du^r}{d\tau} &= -\Gamma^r_{\nu\rho}\, u^\nu\, u^\rho, \label{eq:ode-ur}\\[3pt]
    \frac{du^\theta}{d\tau} &= -\Gamma^\theta_{\nu\rho}\, u^\nu\, u^\rho, \label{eq:ode-utheta}\\[3pt]
    \frac{du^\phi}{d\tau} &= -\Gamma^\phi_{\nu\rho}\, u^\nu\, u^\rho, \label{eq:ode-uphi}\\[3pt]
    \frac{dS^\mu}{d\tau} &= -\Gamma^\mu_{\nu\rho}\, S^\nu\, u^\rho,
    \quad \mu = t, r, \theta, \phi. \label{eq:ode-spin}
\end{align}
Dalam setiap persamaan, sumasi terhadap indeks berulang $\nu$ dan $\rho$
berjalan dari 0 hingga 3 (konvensi sumasi Einstein).
Simbol Christoffel $\Gamma^\mu_{\nu\rho}$ dievaluasi pada posisi sesaat
$(ct, r, \theta, \phi)$ menggunakan ekspresi dari metrik Kerr.


% ============================================================
% ============================================================
%  SECTION 3.6: FORMULASI METODE INTEGRASI NUMERIK (RK4)
% ============================================================
% ============================================================

\section{Formulasi Metode Integrasi Numerik}
\label{sec:metode-numerik}

Sistem 12 ODE orde satu (Persamaan~\ref{eq:ode-ct}--\ref{eq:ode-spin})
bersifat nonlinier karena simbol Christoffel $\Gamma^\mu_{\nu\rho}$
merupakan fungsi rasional dari koordinat $(r, \theta)$.
Sistem ini tidak memiliki solusi analitik umum dan harus diselesaikan
secara numerik.
Pada bagian ini disajikan formulasi metode Runge-Kutta orde empat (RK4)
yang digunakan untuk integrasi sistem tersebut
\parencite{Butcher2016NumericalMethods}.


\subsection{Formulasi Runge-Kutta Orde Empat}
\label{subsec:rk4}

Diberikan masalah nilai awal:
\begin{equation}
    \frac{d\mathbf{Y}}{d\tau} = \mathbf{F}(\mathbf{Y}),
    \qquad \mathbf{Y}(\tau_0) = \mathbf{Y}_0,
    \label{eq:ivp}
\end{equation}
dengan $\mathbf{Y} \in \mathbb{R}^{12}$ dan $\mathbf{F}: \mathbb{R}^{12} \to \mathbb{R}^{12}$.
Metode RK4 mengaproksimasi solusi pada titik diskret
$\tau_n = \tau_0 + n\,\Delta\tau$ melalui iterasi berikut:

\begin{equation}
    \boxed{
    \mathbf{Y}_{n+1} = \mathbf{Y}_n
    + \frac{\Delta\tau}{6}\left(\mathbf{k}_1 + 2\,\mathbf{k}_2
    + 2\,\mathbf{k}_3 + \mathbf{k}_4\right),
    }
    \label{eq:rk4-update}
\end{equation}
dengan empat evaluasi slope:
\begin{align}
    \mathbf{k}_1 &= \mathbf{F}\!\left(\mathbf{Y}_n\right),
    \label{eq:k1} \\[6pt]
    \mathbf{k}_2 &= \mathbf{F}\!\left(\mathbf{Y}_n
        + \frac{\Delta\tau}{2}\,\mathbf{k}_1\right),
    \label{eq:k2} \\[6pt]
    \mathbf{k}_3 &= \mathbf{F}\!\left(\mathbf{Y}_n
        + \frac{\Delta\tau}{2}\,\mathbf{k}_2\right),
    \label{eq:k3} \\[6pt]
    \mathbf{k}_4 &= \mathbf{F}\!\left(\mathbf{Y}_n
        + \Delta\tau\,\mathbf{k}_3\right).
    \label{eq:k4}
\end{align}
Metode ini memiliki orde konvergensi 4, artinya galat lokal
per langkah berskala $\mathcal{O}(\Delta\tau^5)$ dan galat global
berskala $\mathcal{O}(\Delta\tau^4)$.

Setiap evaluasi $\mathbf{F}$ memerlukan perhitungan seluruh
simbol Christoffel $\Gamma^\mu_{\nu\rho}$ pada posisi yang bersesuaian,
sehingga biaya komputasi per langkah adalah $4 \times$
biaya satu evaluasi Christoffel.


\subsection{Langkah Waktu dan Kriteria Kestabilan}
\label{subsec:step-size}

Pemilihan langkah waktu $\Delta\tau$ harus memperhatikan
dua pertimbangan:
\begin{enumerate}
    \item \textbf{Akurasi}: $\Delta\tau$ harus cukup kecil agar galat terpotong
          tidak mendominasi. Secara umum, diperlukan
          $\Delta\tau \ll T_{\text{orbit}}$, di mana $T_{\text{orbit}}$
          adalah periode orbit.
    \item \textbf{Efisiensi}: $\Delta\tau$ tidak boleh terlalu kecil
          agar jumlah langkah total $N = T_{\text{total}}/\Delta\tau$
          tetap dalam batas komputasi yang wajar.
\end{enumerate}

Untuk orbit GP-B dengan $T_{\text{orbit}} \approx 97{,}6$~menit
$\approx 5856$~s, langkah waktu tipikal $\Delta\tau \sim 1$--$10$~s
memberikan $\sim 10^3$--$10^4$ langkah per orbit dan
$\sim 10^6$--$10^7$ langkah untuk waktu integrasi satu tahun.

Sebagai alternatif, dapat digunakan metode adaptif RK45
(Runge-Kutta-Fehlberg) yang tersedia dalam modul
\texttt{scipy.integrate.solve\_ivp} dengan toleransi relatif dan absolut
yang dapat ditentukan pengguna \parencite{SciPy2020}.


\subsection{Verifikasi Kestabilan Numerik}
\label{subsec:verifikasi-kestabilan}

Kestabilan dan akurasi integrasi dimonitor melalui tiga besaran konservatif:
\begin{enumerate}
    \item \textbf{Normalisasi empat-kecepatan}:
    deviasi relatif $\delta_u(\tau) = |g_{\mu\nu} u^\mu u^\nu - c^2|/c^2$
    harus tetap $\ll 1$ sepanjang integrasi
    (Persamaan~\ref{eq:norm-constraint}).

    \item \textbf{Ortogonalitas spin-kecepatan}:
    $\delta_S(\tau) = |g_{\mu\nu} S^\mu u^\nu|$
    harus tetap mendekati nol
    (Persamaan~\ref{eq:ortogonalitas}).

    \item \textbf{Norma spin}:
    $|S|^2 = g_{\mu\nu} S^\mu S^\nu$ harus konstan
    karena transportasi paralel melestarikan norma vektor.
\end{enumerate}

Ketiga kuantitas ini bukan merupakan output fisik yang dicari,
melainkan \textit{diagnostic tool} untuk mendeteksi akumulasi galat numerik.
Hasil pemantauan besaran-besaran ini disajikan pada Bab~\ref{chap:hasil}.


% ============================================================
\section{Parameter Simulasi}
\label{sec:parameter-simulasi}

Parameter yang digunakan dalam simulasi numerik
disajikan pada Tabel~\ref{tab:param-simulasi}.
Nilai-nilai ini dipilih agar konsisten dengan konfigurasi orbit GP-B
dan parameter fisik Bumi yang digunakan dalam literatur.

\begin{table}[ht]
\centering
\caption{Parameter simulasi numerik.}
\label{tab:param-simulasi}
\begin{tabular}{lll}
\toprule
\textbf{Parameter} & \textbf{Nilai} & \textbf{Keterangan} \\
\midrule
Massa Bumi ($M$) & $5{,}972 \times 10^{24}$~kg & --- \\
Momentum sudut ($J$) & $5{,}86 \times 10^{33}$~kg\,m$^2$\,s$^{-1}$ & --- \\
Radius orbit ($r_0$) & $7{,}013 \times 10^6$~m & $R_\oplus + 642$~km \\
Inklinasi orbit & $90°$ (polar) & sesuai GP-B \\
Parameter rotasi ($a$) & $3{,}26$~m & $J/(Mc)$ \\
$r_s/r_0$ & $1{,}26 \times 10^{-9}$ & medan sangat lemah \\
$a/r_0$ & $4{,}65 \times 10^{-7}$ & rotasi sangat lambat \\
Langkah waktu ($\Delta\tau$) & $1$--$10$~s & adaptif atau tetap \\
Waktu integrasi & $3{,}156 \times 10^7$~s & $\approx 1$ tahun \\
Metode integrasi & RK4 / RK45 & via SciPy \\
\bottomrule
\end{tabular}
\end{table}

Kondisi awal untuk posisi dipilih sebagai orbit sirkular
pada bidang ekuator ($\theta_0 = \pi/2$, $\dot{r}_0 = 0$, $\dot{\theta}_0 \neq 0$
untuk orbit polar), dengan kecepatan orbit Keplerian
$v_0 = \sqrt{GM/r_0} \approx 7{,}536 \times 10^3$~m/s.
Kondisi awal vektor spin $S^\mu(0)$ dipilih ortogonal terhadap
$u^\mu(0)$ dan mengarah ke bintang referensi (IM Pegasi),
konsisten dengan konfigurasi eksperimen GP-B
\parencite{Everitt2011GravityProbeB}.
